The function in question is:\newline
%\begin{verbatim}
%\begin{tcolorbox}[colback=lightgray,colframe=black,sharp corners]
%move\_position\_by\_delta\_on\_axis(\newline $\-\-\-$\&self, axis:CameraAxis, delta:f32\newline) -$>$ Result$<$Vector$<$f32$>$, MatrixError$>$
%\end{tcolorbox}
\noindent\fbox{\parbox{\columnwidth}{move\_position\_by\_delta\_on\_axis(\newline $\-\-\-\>\>\-\-\-\>\>$   \&self, axis:CameraAxis, delta:f32\newline) -$>$ Result$<$Vector$<$f32$>$, MatrixError$>$}}\newline
%\end{verbatim}


The function takes inputs of:
\begin{enumerate}
    \item \fbox{\&self}, of type Self, which is type CameraInfoMatrix
    \item \fbox{axis:CameraAxis}, which has enum variants of Forward, Right, and Up
    \item \fbox{delta:f32}, which is the distance delta to move along the specified axis
\end{enumerate}
$\-$\newline
The function produces an output of either:
\begin{enumerate}
    \item \fbox{Ok(Vector$<$f32$>$)}, which is the vector that gets added to the current position to translate it to the new position
    \item \fbox{Err(MatrixError)}, signalling that a math error occured
\end{enumerate}